\documentclass[11pt,letterpaper]{article}
\usepackage[utf8]{inputenc}
\usepackage[T1]{fontenc}
\usepackage[spanish,es-tabla]{babel}
\usepackage{lmodern}
\usepackage[margin=2.2cm]{geometry}
\usepackage{amsmath}
\usepackage{graphicx}
\usepackage{booktabs}
\usepackage{array}
\usepackage{float}
\usepackage{xcolor}
\usepackage{listings}
\usepackage{tikz}
\usetikzlibrary{shapes.geometric,shapes.misc,arrows.meta,positioning,calc}
\usepackage{hyperref}

\definecolor{verde}{RGB}{46,125,90}
\definecolor{gris}{RGB}{90,90,90}

\lstdefinestyle{vhdl}{
  language=VHDL,
  basicstyle=\ttfamily\footnotesize,
  keywordstyle=\color{blue!70!black}\bfseries,
  stringstyle=\color{red!60!black},
  numbers=left,
  numberstyle=\tiny\color{gris},
  frame=single,
  breaklines=true,
  tabsize=4,
  showstringspaces=false
}

\tikzset{
  flecha/.style={-{Stealth[length=2.2mm]}, thick},
  linea/.style={thick},
  estado/.style={draw, thick, minimum width=2.6cm, minimum height=0.85cm, align=center},
  decision/.style={draw, thick, diamond, aspect=1.6, inner sep=1pt, minimum width=1.5cm, minimum height=0.95cm},
  condicional/.style={draw, thick, rounded rectangle, minimum width=1.6cm, minimum height=0.6cm},
  nombre/.style={draw, ellipse, color=verde, font=\footnotesize\bfseries, inner sep=1.5pt},
  codigo/.style={font=\footnotesize\ttfamily},
  rama/.style={font=\footnotesize\bfseries}
}

\newcommand{\nb}[1]{\overline{#1}}

\begin{document}

\begin{titlepage}
\centering
{\Large\textbf{Universidad Nacional Autónoma de México}}\\[0.3cm]
{\large Facultad de Ingeniería}\\[1.5cm]
{\large Organización y Arquitectura de Computadoras}\\[0.3cm]
{\large Laboratorio}\\[2cm]
{\LARGE\textbf{Práctica No. 4}}\\[0.5cm]
{\Large Construcción de Máquinas de Estados usando Memorias:\\[0.2cm] Direccionamiento Entrada-Estado}\\[2.5cm]
\begin{tabular}{rl}
\textbf{Integrantes:} & Domínguez Palacios Jesús Alejandro \\
 & \rule{7cm}{0.4pt} \\
 & \rule{7cm}{0.4pt} \\[0.4cm]
\textbf{Profesor(a):} & \rule{7cm}{0.4pt} \\[0.4cm]
\textbf{Grupo:} & \rule{3cm}{0.4pt} \\[0.4cm]
\textbf{Fecha de entrega:} & 27 de octubre de 2026
\end{tabular}
\vfill
\end{titlepage}

\section{Objetivo}
Familiarizar al alumno en la construcción de máquinas de estados usando direccionamiento de memorias con el método de direccionamiento entrada-estado.

\section{Ejercicio 1: Diseño de la carta ASM}

Se diseñó una carta ASM con las siguientes características:

\begin{itemize}
  \item \textbf{Estados:} 8 (EST0 a EST7).
  \item \textbf{Entradas:} 4 ($A$, $B$, $C$ y $D$), más una variable auxiliar ($aux$) para los estados sin entrada.
  \item \textbf{Salidas:} 4. Dos con lógica estándar ($S_0$ y $S_1$) y dos con lógica negada ($\nb{U_0}$ y $\nb{U_1}$).
  \item \textbf{Salidas condicionales:} 3. $S_1$ en EST0 cuando $B=1$, $S_0$ en EST3 cuando $A=1$ y $\nb{U_1}$ en EST4 cuando $D=0$.
\end{itemize}

Como el método entrada-estado solo permite una entrada por estado, cada estado prueba una sola variable. Los estados EST1, EST5, EST6 y EST7 no prueban ninguna entrada, por lo que en ellos se prueba la variable auxiliar.

\begin{figure}[H]
\centering
\begin{tikzpicture}[scale=0.95, every node/.style={transform shape}]
  \node[estado] (e0) at (0,0) {$S_0$};
  \node[decision] (dB) at (0,-1.6) {$B$};
  \node[condicional] (cS1) at (-3.6,-2.7) {$S_1$};
  \node[estado] (e1) at (-3.6,-4.0) {$S_0,\ \nb{U_1}$};
  \node[estado] (e3) at (3.6,-3.1) {$\nb{U_0}$};
  \node[decision] (dA) at (3.6,-4.6) {$A$};
  \node[condicional] (cS0) at (3.6,-6.0) {$S_0$};
  \node[estado] (e2) at (-3.6,-5.7) {};
  \node[decision] (dC) at (-3.6,-7.2) {$C$};
  \node[estado] (e4) at (3.6,-7.4) {$S_1,\ \nb{U_0}$};
  \node[decision] (dD) at (3.6,-8.9) {$D$};
  \node[estado] (e5) at (-3.6,-8.8) {$S_0,\ S_1$};
  \node[estado] (e6) at (-3.6,-10.4) {$S_1,\ \nb{U_1}$};
  \node[condicional] (cU1) at (3.6,-10.2) {$\nb{U_1}$};
  \node[estado] (e7) at (3.6,-11.5) {$S_0,\ S_1,\ \nb{U_0},\ \nb{U_1}$};

  \foreach \e/\n/\c in {e0/EST0/000, e1/EST1/001, e3/EST3/010, e2/EST2/011, e4/EST4/100, e5/EST5/101, e7/EST7/110, e6/EST6/111}{
    \node[nombre, anchor=south east] at ($(\e.north west)+(0.05,0.05)$) {\n};
    \node[codigo, anchor=south east] at ($(\e.north east)+(0,0.02)$) {\c};
  }

  \draw[flecha] (e0) -- (dB);
  \draw[linea] (dB.west) -- node[rama, above, pos=0.12] {1} (-3.6,-1.6);
  \draw[flecha] (-3.6,-1.6) -- (cS1);
  \draw[flecha] (cS1) -- (e1);
  \draw[linea] (dB.east) -- node[rama, above, pos=0.12] {0} (3.6,-1.6);
  \draw[flecha] (3.6,-1.6) -- (e3);

  \draw[flecha] (e3) -- (dA);
  \draw[flecha] (dA.south) -- node[rama, right, pos=0.3] {1} (cS0);
  \draw[flecha] (cS0) -- (e4);
  \draw[flecha] (dA.west) -- node[rama, above, pos=0.3] {0} ++(-0.6,0) |- (e2.east);

  \draw[flecha] (e1) -- (e2);
  \draw[flecha] (e2) -- (dC);
  \draw[flecha] (dC.south) -- node[rama, right, pos=0.3] {0} (e5);
  \draw[flecha] (dC.east) -- node[rama, above, pos=0.2] {1} (0,-7.2) |- (e4.west);

  \draw[flecha] (e4) -- (dD);
  \draw[flecha] (dD.south) -- node[rama, right, pos=0.3] {0} (cU1);
  \draw[flecha] (cU1) -- (e7);
  \draw[flecha] (dD.west) -- node[rama, above, pos=0.2] {1} (0.6,-8.9) |- (e6.east);

  \draw[flecha] (e5) -- (e6);
  \draw[linea] (e6.west) -- (-6.1,-10.4);
  \draw[linea] (e7.south) -- (3.6,-12.6) -- (-6.1,-12.6) -- (-6.1,-10.4);
  \draw[flecha] (-6.1,-10.4) -- (-6.1,0) -- (e0.west);
\end{tikzpicture}
\caption{Carta ASM de la máquina de estados.}
\label{fig:asm}
\end{figure}

\section{Ejercicio 2: Tamaño y contenido de la memoria}

\subsection{Asignación binaria}

\begin{table}[H]
\centering
\begin{tabular}{ccc}
\toprule
\textbf{Entrada} & \textbf{Código (prueba)} & \textbf{Switch} \\
\midrule
$A$   & 000 & SW0 \\
$B$   & 001 & SW1 \\
$C$   & 010 & SW2 \\
$D$   & 011 & SW3 \\
$aux$ & 100 & Constante 0 \\
\bottomrule
\end{tabular}
\hspace{1.5cm}
\begin{tabular}{cc}
\toprule
\textbf{Estado} & \textbf{Código} \\
\midrule
EST0 & 000 \\
EST1 & 001 \\
EST3 & 010 \\
EST2 & 011 \\
EST4 & 100 \\
EST5 & 101 \\
EST7 & 110 \\
EST6 & 111 \\
\bottomrule
\end{tabular}
\caption{Asignación binaria de entradas y estados.}
\end{table}

\begin{table}[H]
\centering
\begin{tabular}{cccc}
\toprule
\textbf{Salida} & \textbf{Lógica} & \textbf{Activa} & \textbf{Inactiva} \\
\midrule
$S_0$ & Estándar & 1 & 0 \\
$S_1$ & Estándar & 1 & 0 \\
$\nb{U_0}$ & Negada & 0 & 1 \\
$\nb{U_1}$ & Negada & 0 & 1 \\
\bottomrule
\end{tabular}
\caption{Tipo de lógica de las salidas.}
\end{table}

\subsection{Tamaño de la memoria}

Como la carta tiene salidas condicionales, se usa el direccionamiento entrada-estado con salidas condicionadas. La palabra de memoria tiene dos campos de salidas: salidas falsas (cuando la entrada probada vale 0) y salidas verdaderas (cuando vale 1). La dirección de la memoria es únicamente el estado presente.

\begin{itemize}
  \item $n$ = bits de estado $= 3$
  \item $m$ = prueba + liga falsa + liga verdadera + salidas falsas + salidas verdaderas $= 3 + 3 + 3 + 4 + 4 = 17$
\end{itemize}

\[
\text{Memoria} = 2^n \times m = 2^3 \times 17 = 8 \times 17 = 136 \text{ bits}
\]

\subsection{Contenido de la memoria}

En los estados que prueban la variable auxiliar se coloca la misma liga y las mismas salidas en ambos campos, por lo que el valor de $aux$ no altera el funcionamiento.

\begin{table}[H]
\centering
\setlength{\tabcolsep}{5pt}
\begin{tabular}{c|c|c|c|c|cccc|cccc}
\toprule
 & \textbf{Dirección} & & & & \multicolumn{4}{c|}{\textbf{Salidas falsas}} & \multicolumn{4}{c}{\textbf{Salidas verdaderas}} \\
\textbf{Estado} & \textbf{Edo. presente} & \textbf{Prueba} & \textbf{LF} & \textbf{LV} & $S_0$ & $S_1$ & $\nb{U_0}$ & $\nb{U_1}$ & $S_0$ & $S_1$ & $\nb{U_0}$ & $\nb{U_1}$ \\
\midrule
EST0 & 000 & 001 & 010 & 001 & 1 & 0 & 1 & 1 & 1 & 1 & 1 & 1 \\
EST1 & 001 & 100 & 011 & 011 & 1 & 0 & 1 & 0 & 1 & 0 & 1 & 0 \\
EST3 & 010 & 000 & 011 & 100 & 0 & 0 & 0 & 1 & 1 & 0 & 0 & 1 \\
EST2 & 011 & 010 & 101 & 100 & 0 & 0 & 1 & 1 & 0 & 0 & 1 & 1 \\
EST4 & 100 & 011 & 110 & 111 & 0 & 1 & 0 & 0 & 0 & 1 & 0 & 1 \\
EST5 & 101 & 100 & 111 & 111 & 1 & 1 & 1 & 1 & 1 & 1 & 1 & 1 \\
EST7 & 110 & 100 & 000 & 000 & 1 & 1 & 0 & 0 & 1 & 1 & 0 & 0 \\
EST6 & 111 & 100 & 000 & 000 & 0 & 1 & 1 & 0 & 0 & 1 & 1 & 0 \\
\bottomrule
\end{tabular}
\caption{Contenido de la memoria ($8 \times 17$ bits).}
\label{tab:memoria}
\end{table}

\section{Diagrama lógico}

\begin{figure}[H]
\centering
\begin{tikzpicture}[
  bloque/.style={draw, thick, minimum height=0.9cm, align=center, font=\small},
  campo/.style={draw, thick, minimum height=0.8cm, align=center, font=\scriptsize},
  mux/.style={draw, thick, trapezium, trapezium angle=65, shape border rotate=180, minimum height=0.9cm, align=center, font=\scriptsize},
  bus/.style={-{Stealth[length=2mm]}, thick},
  every node/.style={font=\scriptsize}
]
  \node[bloque, minimum width=2.4cm] (reg) at (0,0) {Registro de\\estado (3 bits)};
  \node[bloque, minimum width=9.6cm, anchor=west] (mem) at (2.6,0) {Memoria $8 \times 17$};
  \node[campo, minimum width=1.6cm, anchor=north west] (pr) at (mem.south west) {Prueba\\(3)};
  \node[campo, minimum width=1.9cm, anchor=west] (lf) at (pr.east) {Liga falsa\\(3)};
  \node[campo, minimum width=2.1cm, anchor=west] (lv) at (lf.east) {Liga verdadera\\(3)};
  \node[campo, minimum width=2.0cm, anchor=west] (sf) at (lv.east) {Salidas falsas\\(4)};
  \node[campo, minimum width=2.0cm, anchor=west] (sv) at (sf.east) {Salidas\\verdaderas (4)};
  \draw[bus] (reg.east) -- node[above] {3} (mem.west);

  \node[mux, trapezium angle=78, minimum width=1.9cm, minimum height=1.3cm] (msel) at (3.4,-3.2) {Selector de\\entradas\\MUX 8:1};
  \node[mux, minimum width=2.4cm] (mlig) at (6.15,-3.2) {Selector\\de liga};
  \node[mux, minimum width=2.4cm] (msal) at (10.2,-3.2) {Selector\\de salidas};
  \node[bloque, minimum width=2.4cm] (rsal) at (10.2,-5.1) {Registro de\\salidas};
  \node[align=center] (leds) at (10.2,-6.5) {$S_0$ $S_1$ $\nb{U_0}$ $\nb{U_1}$\\(LEDR3..0)};

  \draw[bus] (pr.south) -- node[right] {3} (msel.north);
  \foreach \t/\y in {A/0.48, B/0.24, C/0, D/-0.24, aux/-0.48}{
    \draw[bus] (1.2,-3.2+\y) node[left] {$\t$} -- (2.3,-3.2+\y);
  }
  \draw[bus] (lf.south) -- ++(0,-0.9) -| (5.55,-2.75);
  \draw[bus] (lv.south) -- ++(0,-0.9) -| (6.75,-2.75);
  \draw[bus] (sf.south) -- ++(0,-0.9) -| (9.6,-2.75);
  \draw[bus] (sv.south) -- ++(0,-0.9) -| (10.8,-2.75);

  \draw[bus] (msel.east) -- node[above, pos=0.4] {$q$} (mlig.west);
  \fill (4.85,-3.2) circle (1.5pt);
  \draw[bus] (4.85,-3.2) -- (4.85,-4.3) -- (8.4,-4.3) |- (msal.west);

  \draw[bus] (msal.south) -- node[right] {4} (rsal.north);
  \draw[bus] (rsal.south) -- (leds.north);

  \draw[bus] (mlig.south) -- (6.15,-6.2) -- node[below, pos=0.5] {Estado siguiente (3)} (0.3,-6.2) -- (0.3,-0.45);

  \node[anchor=east] (clk) at (-1.0,-7.2) {CLOCK\_50 / KEY1};
  \draw[thick] (clk.east) -- (8.7,-7.2);
  \fill (-0.3,-7.2) circle (1.5pt);
  \draw[bus] (-0.3,-7.2) -- (-0.3,-0.45);
  \draw[bus] (8.7,-7.2) |- (rsal.west);

  \node[anchor=east] (rst) at (-1.0,1.2) {KEY0 (reset)};
  \draw[bus] (rst.east) -| (reg.north);
\end{tikzpicture}
\caption{Diagrama lógico del direccionamiento entrada-estado con salidas condicionadas.}
\label{fig:logico}
\end{figure}

El registro de estado direcciona la memoria. El campo de prueba controla el selector de entradas (MUX 8:1), que entrega la entrada probada $q$. Con $q=0$ el selector de liga toma la liga falsa y el selector de salidas toma las salidas falsas; con $q=1$ toma la liga y las salidas verdaderas. La liga elegida se carga en el registro de estado cada vez que se presiona KEY1, y las salidas pasan por un registro antes de llegar a los LEDs.

\section{Ejercicio 3: Implementación en Quartus}

El contenido de la memoria de la Tabla~\ref{tab:memoria} se escribe como una constante en VHDL. Cada palabra está formada por la concatenación de los campos prueba, liga falsa, liga verdadera, salidas falsas y salidas verdaderas.

\begin{lstlisting}[style=vhdl]
library ieee;
use ieee.std_logic_1164.all;
use ieee.numeric_std.all;

entity practica4_asm is
    generic (ANCHO : natural := 20);
    port (
        CLOCK_50 : in  std_logic;
        KEY      : in  std_logic_vector(1 downto 0);
        SW       : in  std_logic_vector(3 downto 0);
        LEDR     : out std_logic_vector(9 downto 0)
    );
end entity;

architecture entrada_estado of practica4_asm is
    type memoria_t is array (0 to 7) of std_logic_vector(16 downto 0);
    constant memoria : memoria_t := (
        "001" & "010" & "001" & "1011" & "1111",
        "100" & "011" & "011" & "1010" & "1010",
        "000" & "011" & "100" & "0001" & "1001",
        "010" & "101" & "100" & "0011" & "0011",
        "011" & "110" & "111" & "0100" & "0101",
        "100" & "111" & "111" & "1111" & "1111",
        "100" & "000" & "000" & "1100" & "1100",
        "100" & "000" & "000" & "0110" & "0110"
    );
    constant aux : std_logic := '0';
    signal estado_presente    : std_logic_vector(2 downto 0) := "000";
    signal palabra            : std_logic_vector(16 downto 0);
    signal prueba             : std_logic_vector(2 downto 0);
    signal liga_falsa         : std_logic_vector(2 downto 0);
    signal liga_verdadera     : std_logic_vector(2 downto 0);
    signal liga               : std_logic_vector(2 downto 0);
    signal salidas_falsas     : std_logic_vector(3 downto 0);
    signal salidas_verdaderas : std_logic_vector(3 downto 0);
    signal salidas_sel        : std_logic_vector(3 downto 0);
    signal salidas            : std_logic_vector(3 downto 0) := "1011";
    signal a, b, c, d, q      : std_logic;
    signal s     : std_logic_vector(1 downto 0) := "11";
    signal ant   : std_logic := '1';
    signal pulso : std_logic := '0';
    signal cont  : unsigned(ANCHO-1 downto 0) := (others => '0');
begin
    a <= SW(0);
    b <= SW(1);
    c <= SW(2);
    d <= SW(3);

    palabra            <= memoria(to_integer(unsigned(estado_presente)));
    prueba             <= palabra(16 downto 14);
    liga_falsa         <= palabra(13 downto 11);
    liga_verdadera     <= palabra(10 downto 8);
    salidas_falsas     <= palabra(7 downto 4);
    salidas_verdaderas <= palabra(3 downto 0);

    with prueba select
        q <= a   when "000",
             b   when "001",
             c   when "010",
             d   when "011",
             aux when others;

    liga        <= liga_verdadera     when q = '1' else liga_falsa;
    salidas_sel <= salidas_verdaderas when q = '1' else salidas_falsas;

    process (CLOCK_50)
    begin
        if rising_edge(CLOCK_50) then
            s     <= s(0) & KEY(1);
            cont  <= cont + 1;
            pulso <= '0';
            if cont = 0 then
                ant <= s(1);
                if ant = '1' and s(1) = '0' then
                    pulso <= '1';
                end if;
            end if;
        end if;
    end process;

    process (CLOCK_50, KEY)
    begin
        if KEY(0) = '0' then
            estado_presente <= "000";
        elsif rising_edge(CLOCK_50) then
            if pulso = '1' then
                estado_presente <= liga;
            end if;
        end if;
    end process;

    process (CLOCK_50)
    begin
        if rising_edge(CLOCK_50) then
            salidas <= salidas_sel;
        end if;
    end process;

    LEDR(3 downto 0) <= salidas;
    LEDR(6 downto 4) <= "000";
    LEDR(9 downto 7) <= estado_presente;
end architecture;
\end{lstlisting}

\subsection{Asignación de pines (DE10-Lite)}

\begin{table}[H]
\centering
\begin{tabular}{lll@{\hspace{1.5cm}}lll}
\toprule
\textbf{Señal} & \textbf{Pin} & \textbf{Función} & \textbf{Señal} & \textbf{Pin} & \textbf{Función} \\
\midrule
CLOCK\_50 & PIN\_P11 & Reloj 50 MHz & LEDR0 & PIN\_A8  & $\nb{U_1}$ \\
KEY0      & PIN\_B8  & Reset        & LEDR1 & PIN\_A9  & $\nb{U_0}$ \\
KEY1      & PIN\_A7  & Avanzar      & LEDR2 & PIN\_A10 & $S_1$ \\
SW0       & PIN\_C10 & $A$          & LEDR3 & PIN\_B10 & $S_0$ \\
SW1       & PIN\_C11 & $B$          & LEDR4 & PIN\_D13 & 0 \\
SW2       & PIN\_D12 & $C$          & LEDR5 & PIN\_C13 & 0 \\
SW3       & PIN\_C12 & $D$          & LEDR6 & PIN\_E14 & 0 \\
          &          &              & LEDR7 & PIN\_D14 & Estado bit 0 \\
          &          &              & LEDR8 & PIN\_A15 & Estado bit 1 \\
          &          &              & LEDR9 & PIN\_B11 & Estado bit 2 \\
\bottomrule
\end{tabular}
\caption{Asignación de pines en la tarjeta DE10-Lite.}
\end{table}

\section{Ejercicio 4: Prueba en la FPGA}

Se cargó el diseño en la tarjeta DE10-Lite. Con todos los switches en 0 se realizó la siguiente secuencia, que recorre los ocho estados y las tres salidas condicionales.

\begin{table}[H]
\centering
\begin{tabular}{clcc}
\toprule
\textbf{Paso} & \textbf{Acción} & \textbf{Estado (LEDR9..7)} & \textbf{$S_0\,S_1\,\nb{U_0}\,\nb{U_1}$ (LEDR3..0)} \\
\midrule
1  & Reset (KEY0), $B=0$            & EST0 (000) & 1011 \\
2  & Subir $B$ (salida condicional)  & EST0 (000) & 1111 \\
3  & Bajar $B$, presionar KEY1       & EST3 (010) & 0001 \\
4  & Subir $A$ (salida condicional)  & EST3 (010) & 1001 \\
5  & Bajar $A$, presionar KEY1       & EST2 (011) & 0011 \\
6  & $C=0$, presionar KEY1           & EST5 (101) & 1111 \\
7  & Presionar KEY1                  & EST6 (111) & 0110 \\
8  & Presionar KEY1                  & EST0 (000) & 1011 \\
9  & $B=1$, presionar KEY1           & EST1 (001) & 1010 \\
10 & Presionar KEY1                  & EST2 (011) & 0011 \\
11 & $C=1$, presionar KEY1           & EST4 (100) & 0100 \\
12 & Subir $D$ (salida condicional)  & EST4 (100) & 0101 \\
13 & Bajar $D$, presionar KEY1       & EST7 (110) & 1100 \\
14 & Presionar KEY1                  & EST0 (000) & 1011 \\
\bottomrule
\end{tabular}
\caption{Secuencia de prueba en la tarjeta.}
\end{table}

Los caminos restantes se comprueban de la misma forma: desde EST3 con $A=1$ se pasa a EST4, y desde EST4 con $D=1$ se pasa a EST6.

\section{Conclusiones}

% Escribir aquí las conclusiones de cada integrante.

\end{document}